% Review tiga paper rujukan proyek: Audebert (2020), Kanchi (2022), Kashyap (2026).
% Memakai gaya template FILKOM UB yang sama dengan ../../proposal/analisis-paper.
% Kompilasi: latexmk main.tex   (pdfLaTeX + biber)
\documentclass[12pt,a4paper,oneside]{report}
\usepackage{gaya/ubtesis}
\usepackage{makecell}
\graphicspath{{../gambar/}}
\addbibresource{referensi.bib}

\renewcommand\thesection{\arabic{section}}
\counterwithout{table}{chapter}
\counterwithout{figure}{chapter}
\counterwithout{equation}{chapter}
\renewcommand\theequation{\arabic{equation}}
\defbibheading{laporan}{\section*{Daftar Referensi}}
\pretocmd{\section}{\needspace{12\baselineskip}}{}{}

% ------------------------------------------------------------- data sampul
\newcommand\judul{Review Paper: Klasifikasi Dokumen Multimodal Teks OCR dan Citra\\[4pt]
  \normalsize Audebert (2020), Kanchi (2022), dan Kashyap (2026)}
\newcommand\matakuliah{ANALISIS \& PEMROSESAN DATA MULTIMODAL}
\newcommand\pengampu{Dr. Eng. Fitra Abdurrachman Bachtiar, S.T., M.Eng.}
\newcommand\programstudi{PROGRAM STUDI MAGISTER ILMU KOMPUTER}
\newcommand\departemen{DEPARTEMEN TEKNIK INFORMATIKA}
\newcommand\fakultas{FAKULTAS ILMU KOMPUTER}
\newcommand\universitas{UNIVERSITAS BRAWIJAYA}
\newcommand\anggota{%
  Yulian Anggreaif Rudamaga & 256150100011001 \\
  Zalzabila Rani             & 256150100011008 \\
  Ramadhani Al-Qadri         & 256150100011010 \\}

\begin{document}
\margininduk
% Sampul, mengikuti susunan sampul Progress 1 kelompok.
\begin{titlepage}
\thispagestyle{empty}
\setlength{\parskip}{0pt}
\centering
\vspace*{0.4cm}
{\bfseries\ukuran{14}\matakuliah\par}
\vspace{0.8cm}
{\bfseries\ukuran{16}\judul\par}
\vspace{1.8cm}
\includegraphics[width=4.5cm]{logo-ub}\par
\vspace{1.8cm}
{\ukuran{12}Disusun oleh:\par\vspace{0.25cm}
 \begin{tabular}{@{}l@{\quad:\quad}l@{}}\anggota\end{tabular}\par}
\vspace{1.2cm}
{\ukuran{12}Dosen Pengampu:\par
 \pengampu\par}
\vfill
{\bfseries\ukuran{14}\programstudi\par
 \departemen\par
 \fakultas\par
 \universitas\par}
\vspace{0.5cm}
\end{titlepage}
\setcounter{page}{1}


\marginisi
\pagenumbering{arabic}
% Review dua paper. Setiap paper dibahas dengan sembilan poin wajib dari dosen, ditambah
% empat poin rekomendasi (kontribusi, detail pelatihan, kelebihan, catatan kritis).
% Semua angka disalin dari tabel/teks paper; tidak ada angka dari eksperimen kelompok.

\section{Pendahuluan}\label{sec:pendahuluan}

Laporan ini mereview dua paper tentang klasifikasi citra dokumen yang memadukan dua
modalitas: citra halaman dan teks hasil OCR. Paper pertama adalah dasar replikasi proyek
pada Fase~1, sedangkan paper kedua adalah referensi utama Fase~2. Identitas keduanya
ada pada Tabel~\ref{tab:identitas}.

\begin{table}[H]
  \centering
  \caption{Identitas paper yang direview}
  \label{tab:identitas}
  \fontsize{10}{12}\selectfont
  \begin{tabularx}{\textwidth}{|l|>{\raggedright\arraybackslash}X|>{\raggedright\arraybackslash}X|}
    \hline
    & \textbf{Paper 1} & \textbf{Paper 2} \\ \hline
    Judul & Multimodal Deep Networks for Text and Image-Based Document Classification
          & Multimodal Fusion and Supervised Contrastive Learning for Robust Document Image Classification \\ \hline
    Penulis & Audebert, Herold, Slimani, Vidal (Quicksign, Prancis)
            & Kashyap, Mittal, Singh (Central University of Punjab, India) \\ \hline
    Venue & ECML PKDD 2019 Workshops, CCIS 1167, Springer, hlm.\ 427--443
          & International Journal of Intelligent Systems (Wiley), artikel 4241437, open access \\ \hline
    Tahun & 2020 (workshop 2019) & 2026 (diterima 13 Juli 2026) \\ \hline
    DOI & 10.1007/978-3-030-43823-4\_35 & 10.1155/int/4241437 \\ \hline
    Sitasi & \textcite{audebert2020} & \textcite{kashyap2026} \\ \hline
  \end{tabularx}
\end{table}

Setiap paper dibahas dengan sembilan poin wajib:
\begin{enumerate*}[label=(\arabic*)]
  \item latar belakang dan permasalahan,
  \item dataset,
  \item modalitas,
  \item pipeline keseluruhan,
  \item model per unimodal,
  \item metode fusi,
  \item metrik evaluasi,
  \item hasil keseluruhan beserta hasil per unimodal, dan
  \item \emph{future work}.
\end{enumerate*}
Kami menambahkan empat poin yang menurut kami perlu untuk review yang utuh:
\begin{enumerate*}[label=(\arabic*), start=10]
  \item kontribusi utama,
  \item detail pelatihan dan lingkungan eksperimen,
  \item kelebihan, dan
  \item catatan kritis (kelemahan, inkonsistensi, ancaman validitas).
\end{enumerate*}
Bagian~\ref{sec:banding} membandingkan kedua paper, dan Bagian~\ref{sec:relevansi}
menjelaskan apa yang kami ambil untuk proyek.

% =====================================================================
\section{Paper 1: Audebert et al.\ (2020)}\label{sec:audebert}

\subsection{Latar Belakang dan Permasalahan}
Banyak proses administrasi (pembukaan rekening, pajak, verifikasi identitas) kini
menerima dokumen yang dipindai atau difoto pengguna. Dokumen ini tidak punya lapisan
teks, sehingga pengecekan dan pengarsipannya masih dikerjakan operator manusia.
Langkah pertama yang perlu diotomasi adalah mengenali jenis dokumen.

Masalah yang diangkat: klasifikasi berbasis tampilan visual saja cukup untuk kelas
yang kasar (paspor vs.\ data bank), tetapi tidak cukup untuk kelas yang halus. Dua
formulir pajak bisa berbagi tata letak dan logo yang sama, sedangkan isinya berbeda.
Informasi pembedanya ada di teks, padahal teks itu tidak tersedia dalam bentuk digital
dan harus diambil lewat OCR yang hasilnya berderau. Pertanyaan penelitiannya: apakah
teks OCR yang berderau tetap bisa menaikkan akurasi klasifikasi citra dokumen jika
digabung dengan fitur visual, dengan biaya komputasi yang tetap rendah?

\subsection{Dataset}
\begin{itemize}
  \item \textbf{Tobacco3482} \parencite{kumar2014}: 3.482 dokumen hitam-putih, 10 kelas
        (\emph{advertisement, email, form, letter, memo, news, note, report, resume,
        scientific}). Protokol: 800 dokumen untuk latih, sisanya untuk uji; hasil dirata-rata
        dari 3 kali percobaan.
  \item \textbf{RVL-CDIP} \parencite{harley2015}: 400.000 dokumen grayscale, 16 kelas,
        25.000 per kelas. Pembagian standar 320.000 latih, 40.000 validasi, 40.000 uji.
  \item \textbf{QS-OCR}: teks hasil Tesseract untuk kedua dataset di atas, dirilis penulis
        di GitHub \texttt{Quicksign/ocrized-text-dataset}. Dataset inilah yang kami pakai
        sebagai modalitas teks (QS-OCR-Small untuk Tobacco3482).
\end{itemize}

\subsection{Modalitas}
Dua modalitas yang keduanya berasal dari satu citra halaman:
\begin{enumerate*}[label=(\alph*)]
  \item \textbf{citra} grayscale halaman dokumen, dan
  \item \textbf{teks} yang diekstraksi dari citra yang sama dengan Tesseract OCR.
\end{enumerate*}
Teks bukan modalitas yang independen, melainkan turunan dari citra, sehingga kualitasnya
dibatasi oleh kualitas OCR. Tata letak (posisi kata) tidak dipakai.

\subsection{Pipeline Keseluruhan}
Gambar~\ref{fig:pipeline-audebert} merangkum pipeline. Citra diproses oleh MobileNetV2,
teks OCR diubah menjadi barisan vektor kata FastText lalu diproses CNN 1D. Kedua vektor
fitur 128 dimensi digabung (konkatenasi) dan diklasifikasi oleh MLP. Seluruh jaringan
dapat dilatih \emph{end-to-end}.

\begin{figure}[H]
  \centering
  % Pipeline Audebert et al. (2020), digambar ulang dari Gambar 1 dan Bagian 3-4 paper.
% Dipakai bersama oleh laporan dan presentasi; butuh TikZ (positioning, arrows.meta).
\resizebox{\linewidth}{!}{%
\begin{tikzpicture}[
  node distance=4mm and 7mm,
  kotak/.style={draw, rounded corners=2pt, align=center, font=\small,
                minimum height=10mm, inner sep=3pt},
  citra/.style={kotak, fill=blue!10},
  teks/.style={kotak, fill=orange!15},
  fusi/.style={kotak, fill=green!15},
  panah/.style={-{Stealth[length=2mm]}, thick}]
  \node[kotak] (in) {Citra dokumen\\(grayscale)};
  % alur citra
  \node[citra, above right=4mm and 9mm of in] (rs) {Resize 384$\times$384\\gray$\to$3 kanal};
  \node[citra, right=of rs] (mb) {MobileNetV2\\(ImageNet, fine-tune)};
  \node[citra, right=of mb] (gap) {GAP 1280-d\\$\to$ FC 128-d};
  % alur teks
  \node[teks, below right=4mm and 9mm of in] (ocr) {Tesseract 4.0\\(LSTM, EN)};
  \node[teks, right=of ocr] (ft) {SpaCy token\\+ FastText (Common Crawl)};
  \node[teks, right=of ft] (cnn) {CNN 1D 4 lapis\\$\to$ 128-d};
  % fusi
  \node[fusi, right=12mm of $(gap)!0.5!(cnn)$] (cat) {Konkatenasi\\256-d};
  \node[fusi, right=of cat] (out) {MLP +\\softmax};
  \draw[panah] (in) |- (rs);  \draw[panah] (rs) -- (mb);  \draw[panah] (mb) -- (gap);
  \draw[panah] (in) |- (ocr); \draw[panah] (ocr) -- (ft); \draw[panah] (ft) -- (cnn);
  \draw[panah] (gap) -| (cat); \draw[panah] (cnn) -| (cat); \draw[panah] (cat) -- (out);
  \node[font=\footnotesize\itshape, below=1mm of out, align=center] {dilatih\\end-to-end};
\end{tikzpicture}}

  \caption{Pipeline Audebert et al.\ (digambar ulang dari Gambar~1 dan Bagian~3--4 paper)}
  \label{fig:pipeline-audebert}
\end{figure}

\subsection{Model per Unimodal}
\paragraph{Citra.} MobileNetV2 dipilih karena ringan (sekitar setara VGG-16 dalam
akurasi, tetapi jauh lebih cepat). Citra di-\emph{resize} ke $384\times384$ (rasio
aspek dikorbankan), kanal grayscale diduplikasi tiga kali agar bobot ImageNet bisa
dipakai, lalu jaringan di-\emph{fine-tune}. Augmentasi (artefak JPEG, saturasi, kontras)
tidak membantu, sehingga tidak dipakai.

\paragraph{Teks.} Tesseract 4.0 (mesin LSTM, bahasa Inggris, segmentasi halaman penuh).
Penulis menunjukkan derau OCR: rata-rata 136 kata ($\geq$4 huruf) per dokumen, hanya 118
ada di kosakata GloVe dan 114 di kamus Enchant; sekitar 26\,\% korpus tidak ada di kamus
bahasa Inggris. Karena itu dipilih embedding berbasis subkata, \textbf{FastText} (Common
Crawl), yang tetap memberi vektor masuk akal untuk kata salah eja. Tokenisasi memakai
SpaCy. Dua klasifikator dibandingkan:
\begin{itemize}
  \item MLP atas embedding dokumen (rata-rata berbobot ala Arora + PCA): 2048 neuron per
        lapis, ReLU, Dropout, BatchNorm, lapis akhir 128-d;
  \item \textbf{CNN 1D} atas barisan kata (dipilih): 4 lapis, kernel 12, 512 kanal, ReLU,
        max-pooling stride 2, lalu max-pooling sepanjang waktu, Dropout, dan FC. Input
        dipotong/di-\emph{padding} ke 500 kata.
\end{itemize}

\subsection{Metode Fusi}
Fusi tingkat fitur (\emph{intermediate/feature-level}). Vektor teks 128-d (lapis
sebelum softmax CNN 1D) dan vektor citra 1280-d hasil \emph{global average pooling} yang
diproyeksikan ke 128-d oleh FC, lalu \textbf{dikonkatenasi} dan diumpankan ke MLP.
Penulis juga mencoba \emph{adaptive averaging} (penjumlahan fitur); pada eksperimen awal
cara ini kalah dari model citra saja, dan penulis menduga ruang fitur teks dan citra
tidak bisa diselaraskan tanpa merusak daya pembedanya.

\subsection{Metrik Evaluasi}
\begin{itemize}
  \item \emph{Overall accuracy} (OA);
  \item F1 seimbang antarkelas (\emph{class-balanced} F1, yaitu macro-F1) dan F1 per kelas
        untuk Tobacco3482;
  \item \textbf{Oracle}: akurasi jika selalu dipilih prediksi yang benar dari salah satu
        model unimodal, sebagai batas atas teoretis fusi;
  \item Waktu proses (latensi per dokumen di CPU dan GPU).
\end{itemize}

\subsection{Hasil Keseluruhan}
Tabel~\ref{tab:aud-tob} dan~\ref{tab:aud-rvl} menyalin hasil paper. Pada kedua dataset
pola hasilnya sama: teks paling lemah, citra lebih kuat, fusi terbaik. Fusi menaikkan OA
Tobacco3482 sebesar 3,3 poin di atas model citra (dari batas atas oracle 7,6 poin), dan
RVL-CDIP sebesar 1,5 poin. Pada F1 per kelas, fusi hampir selalu menyamai atau melebihi
unimodal terbaik; penurunan kecil hanya pada \emph{advertisement} (0,93 vs.\ 0,94).
Kelas yang paling diuntungkan teks adalah \emph{resume} (teks 0,97 vs.\ citra 0,80).

\begin{table}[H]
  \centering
  \caption{Hasil Audebert et al.\ pada Tobacco3482 (Tabel~3 paper): OA (\%), macro-F1, dan F1 per kelas}
  \label{tab:aud-tob}
  \fontsize{9}{11}\selectfont
  \setlength\tabcolsep{3pt}
  \begin{tabular}{|l|c|c|c|c|c|c|c|c|c|c|c|c|}
    \hline
    \textbf{Model} & \textbf{OA} & \textbf{F1} & Adv & Email & Form & Letter & Memo & News & Note & Report & Res & Sci \\ \hline
    CNNs (Harley) & 79,9 & -- & \multicolumn{10}{c|}{--} \\ \hline
    Teks   & 73,8 & 0,71 & 0,60 & 0,96 & 0,76 & 0,71 & 0,79 & 0,67 & 0,62 & 0,43 & 0,97 & 0,57 \\ \hline
    Citra  & 84,5 & 0,82 & 0,94 & 0,96 & 0,85 & 0,83 & 0,90 & 0,89 & 0,83 & 0,61 & 0,80 & 0,62 \\ \hline
    \textbf{Fusi} & \textbf{87,8} & \textbf{0,86} & 0,93 & 0,98 & 0,88 & 0,86 & 0,90 & 0,90 & 0,85 & 0,71 & 0,96 & 0,68 \\ \hline
    Oracle & 92,1 & 0,91 & 0,94 & 0,99 & 0,94 & 0,92 & 0,93 & 0,93 & 0,89 & 0,81 & 0,97 & 0,79 \\ \hline
  \end{tabular}
\end{table}

\begin{table}[H]
  \centering
  \caption{Hasil Audebert et al.\ pada RVL-CDIP (Tabel~2 paper): OA (\%)}
  \label{tab:aud-rvl}
  \fontsize{10}{12}\selectfont
  \begin{tabular}{|c|c|c|c|c|c|}
    \hline
    \textbf{Citra} & \textbf{Teks} & \textbf{Fusi} & CNNs (Harley) & VGG-16 (Afzal) & AlexNet+SPP (Tensmeyer) \\ \hline
    89,1 & 74,6 & \textbf{90,6} & 89,8 & 90,97 & 90,94 \\ \hline
  \end{tabular}
\end{table}

Eksperimen pendahuluan pada Tobacco3482 (Tabel~1 paper): MLP dokumen 70,8\,\% / F1 0,69,
CNN 1D 73,9\,\% / 0,71; MobileNetV2 84,5\,\% / 0,82, dengan augmentasi 83,9\,\% / 0,82.

\paragraph{Waktu proses.} Di CPU Intel i7-8550U, Tesseract butuh $\approx$910~ms per
halaman, inferensi model fusi $\approx$360~ms (130~ms teks, 230~ms citra), total
$\approx$1,3~detik. Di GPU Titan X model fusi berjalan 110~ms. Sebagai pembanding,
Xception butuh 630~ms dan LSTM setara butuh beberapa detik. Artinya, hambatan utama ada
di OCR, bukan di jaringan.

\subsection{Future Work}
\begin{itemize}
  \item Uji pada data dunia nyata (foto ponsel, rotasi, pencahayaan beragam); augmentasi
        terutama rotasi menjadi wajib di sana.
  \item Praproses tata letak/segmentasi halaman (mis.\ dhSegment) untuk menormalkan
        orientasi dan memotong dokumen sebelum klasifikasi.
  \item Koreksi ejaan parsial pasca-OCR (SymSpell) yang dikombinasikan dengan NER SpaCy agar
        nama diri, singkatan, dan kode tidak ikut "dikoreksi".
  \item Embedding yang khusus domain, berbasis karakter, atau tahan derau OCR.
  \item Mengganti ekstraktor fitur (CNN lain, model NLP subkata yang dapat di-\emph{fine-tune}),
        karena arsitekturnya modular.
\end{itemize}

\subsection{Kontribusi Utama}
\begin{enumerate*}[label=(\alph*)]
  \item baseline multimodal yang pragmatis dari komponen siap pakai (MobileNetV2 +
        FastText) yang dapat dilatih \emph{end-to-end};
  \item bukti bahwa teks OCR berderau tetap menambah akurasi pada dataset kecil maupun
        besar;
  \item rilis dataset teks QS-OCR untuk Tobacco3482 dan RVL-CDIP;
  \item analisis latensi yang menunjukkan sistem layak untuk aplikasi hampir \emph{real-time}
        di CPU.
\end{enumerate*}

\subsection{Detail Pelatihan dan Lingkungan}
TensorFlow 1.12 (Keras), GPU NVIDIA Titan X. Semua model dioptimasi dengan SGD momentum
0,9, laju belajar 0,01, \emph{batch} 40; teks 100 epoch, citra dan fusi 200 epoch.
Hiperparameter disetel manual pada subset Tobacco3482 lalu dikunci untuk semua eksperimen.

\subsection{Kelebihan}
\begin{itemize}
  \item Desain sederhana dan terukur: tiap komponen dipilih dengan eksperimen pendahuluan
        (MLP vs.\ CNN 1D, dengan/tanpa augmentasi, konkatenasi vs.\ penjumlahan).
  \item Melaporkan oracle, sehingga besar kenaikan fusi dapat dibaca relatif terhadap batas
        atasnya, bukan hanya relatif terhadap baseline.
  \item Melaporkan F1 per kelas untuk ketiga model, sehingga komplementaritas modalitas
        terlihat jelas (mis.\ \emph{resume} sangat terbantu teks).
  \item Analisis derau OCR yang kuantitatif (persentase kata di luar kosakata) dan
        pengukuran latensi yang memasukkan OCR.
  \item Data teks dirilis terbuka sehingga hasilnya dapat direplikasi.
\end{itemize}

\subsection{Catatan Kritis}
\begin{itemize}
  \item \textbf{Variansi tidak dilaporkan.} Tobacco3482 dirata-rata dari 3 kali percobaan
        tetapi tanpa simpangan baku, sehingga tidak terlihat apakah kenaikan 3,3 poin
        lebih besar dari variasi antar-\emph{split}. Tidak ada uji signifikansi untuk klaim kenaikan
        "signifikan".
  \item \textbf{Protokol Tobacco3482 kurang rinci.} Disebut "k-fold dengan 800 dokumen
        latih", tetapi cara memilih 800 dokumen (acak/stratifikasi, berapa per kelas) dan
        ada tidaknya set validasi tidak dijelaskan.
  \item \textbf{Inkonsistensi kecil}: OA teks 73,9\,\% di Tabel~1 vs.\ 73,8\,\% di Tabel~3.
  \item \textbf{RVL-CDIP hanya OA}; F1 dan hasil per kelas tidak dilaporkan.
  \item \textbf{Fusi tunggal.} Hanya konkatenasi yang dilaporkan penuh; tidak ada
        perbandingan dengan \emph{late fusion} (rata-rata probabilitas), padahal oracle
        menunjukkan masih ada celah 4,3 poin.
  \item \textbf{OCR tidak divariasikan.} Hanya Tesseract dengan satu konfigurasi; pengaruh
        mesin OCR terhadap hasil fusi tidak diuji. Ini celah yang diisi Fase~2 proyek kami.
  \item Fusi diklaim dilatih \emph{end-to-end}, tetapi FastText tidak di-\emph{fine-tune},
        sehingga cabang teks hanya melatih CNN 1D di atas embedding beku.
\end{itemize}

% =====================================================================
\section{Paper 2: Kashyap et al.\ (2026)}\label{sec:kashyap}

\subsection{Latar Belakang dan Permasalahan}
Volume dokumen digital terus naik dan klasifikasinya perlu diotomasi. Penulis
mengidentifikasi tiga masalah:
\begin{enumerate*}[label=(\alph*)]
  \item pendekatan unimodal tidak tangguh: tampilan visual dokumen sekelas bisa berbeda dan
        dokumen beda kelas bisa mirip, sementara pendekatan teks rentan terhadap kata
        \emph{out-of-vocabulary} (OOV) akibat OCR buruk dan mengabaikan struktur halaman;
  \item model besar \emph{vision-language} (LayoutLMv3, DocFormer, DocLLM) akurat tetapi
        berat sehingga sulit dipasang di perangkat \emph{edge}, ponsel, atau IoMT;
  \item fusi awal/menengah rentan terhadap propagasi galat antarmodal jika teks OCR
        berderau.
\end{enumerate*}
Celah yang diangkat: literatur fokus memaksimalkan akurasi dengan arsitektur yang makin
kompleks, dan jarang membahas \emph{trade-off} akurasi--efisiensi--sumber daya. Penulis
menawarkan dua kerangka terpisah agar praktisi bisa memilih sesuai kebutuhan: satu
mengejar akurasi, satu mengejar efisiensi.

\subsection{Dataset}
\begin{table}[H]
  \centering
  \caption{Dataset Kashyap et al.\ (Tabel~2 paper)}
  \label{tab:kas-data}
  \fontsize{10}{12}\selectfont
  \begin{tabular}{|l|l|r|r|r|r|c|}
    \hline
    \textbf{Dataset} & \textbf{Jenis} & \textbf{Total} & \textbf{Latih} & \textbf{Val} & \textbf{Uji} & \textbf{Kelas} \\ \hline
    RVL-CDIP    & sekunder & 399.743 & 319.754 & 39.995 & 39.995 & 16 \\ \hline
    Tobacco3482 & sekunder & 3.482   & --      & --     & 3.482  & 10 \\ \hline
    RealDoc-16  & primer   & 160     & --      & --     & 160    & 16 \\ \hline
  \end{tabular}
\end{table}
Semua model dilatih \emph{hanya} pada RVL-CDIP. Tobacco3482 dipakai \emph{seluruhnya}
sebagai set uji (bukan protokol 800 latih yang lazim) untuk menguji generalisasi pada
kondisi tidak seimbang dan sumber daya rendah. RealDoc-16 adalah dataset buatan penulis
(10 dokumen per kelas, taksonomi kelas RVL-CDIP) untuk menguji generalisasi lintas
domain; data ini hanya tersedia atas permintaan.

\subsection{Modalitas}
Sama dengan Paper~1: \textbf{citra} halaman dan \textbf{teks} hasil OCR dari citra yang
sama. Pendekatan~1 memakai keduanya; pendekatan~2 (SCSR) hanya memakai teks. Tata letak
juga tidak dipakai secara eksplisit.

\subsection{Pipeline Keseluruhan}
Gambar~\ref{fig:pipeline-kashyap} menunjukkan dua pipeline yang berdiri sendiri.
\begin{itemize}
  \item \textbf{Pendekatan 1 -- multimodal late fusion.} Dua aliran dilatih terpisah:
        EfficientNetB3 untuk citra dan BERT\textsubscript{base} untuk teks OCR, keduanya
        di-\emph{fine-tune} pada RVL-CDIP. Logit kedua aliran dirata-rata, lalu argmax.
  \item \textbf{Pendekatan 2 -- SCSR} (\emph{supervised contrastive semantic retrieval}).
        Model embedding bge-small-en di-\emph{fine-tune} secara kontrastif terawasi dengan
        pasangan anchor--positif sekelas. Teks latih dipotong menjadi \emph{chunk},
        di-embed, dan disimpan di indeks FAISS beserta labelnya. Dokumen uji di-embed,
        10 \emph{chunk} paling mirip diambil, dan labelnya ditentukan oleh voting mayoritas
        (pada dasarnya $k$-NN dengan $k=10$ di ruang embedding hasil pelatihan kontrastif).
\end{itemize}

\begin{figure}[H]
  \centering
  % Dua pipeline Kashyap et al. (2026), digambar ulang dari Gambar 1 dan Gambar 3 paper.
% Dipakai bersama oleh laporan dan presentasi; butuh TikZ (positioning, arrows.meta).
\resizebox{\linewidth}{!}{%
\begin{tikzpicture}[
  node distance=4mm and 6mm,
  kotak/.style={draw, rounded corners=2pt, align=center, font=\small,
                minimum height=10mm, inner sep=3pt},
  citra/.style={kotak, fill=blue!10},
  teks/.style={kotak, fill=orange!15},
  fusi/.style={kotak, fill=green!15},
  ret/.style={kotak, fill=violet!12},
  panah/.style={-{Stealth[length=2mm]}, thick}]
  % ---- Pendekatan 1: late fusion
  \node[font=\small\bfseries, anchor=west] (j1) at (0,1.9) {Pendekatan 1 -- multimodal late fusion};
  \node[kotak] (in) at (0.9,0) {Citra\\dokumen};
  \node[citra, above right=1mm and 8mm of in] (pre) {Resize 300$\times$300\\norm.\ [0,1], RGB};
  \node[citra, right=of pre] (eff) {EfficientNetB3\\fine-tune RVL-CDIP};
  \node[citra, right=of eff] (lv) {logit visual\\16 kelas};
  \node[teks, below right=1mm and 8mm of in] (ocr) {OCR komersial\\+ praproses teks};
  \node[teks, right=of ocr] (bert) {BERT\textsubscript{base}\\fine-tune RVL-CDIP};
  \node[teks, right=of bert] (lt) {logit teks\\16 kelas};
  \node[fusi, right=10mm of $(lv)!0.5!(lt)$] (avg) {Rata-rata logit\\$(z_v+z_t)/2$};
  \node[fusi, right=of avg] (y1) {argmax};
  \draw[panah] (in) |- (pre); \draw[panah] (pre) -- (eff); \draw[panah] (eff) -- (lv);
  \draw[panah] (in) |- (ocr); \draw[panah] (ocr) -- (bert); \draw[panah] (bert) -- (lt);
  \draw[panah] (lv) -| (avg); \draw[panah] (lt) -| (avg); \draw[panah] (avg) -- (y1);
  % ---- Pendekatan 2: SCSR
  \node[font=\small\bfseries, anchor=west] (j2) at (0,-2.4) {Pendekatan 2 -- SCSR (teks saja)};
  \node[teks] (ocr2) at (1.2,-3.6) {Teks OCR\\RVL-CDIP train};
  \node[ret, right=of ocr2] (pair) {319.625 pasangan\\anchor--positif};
  \node[ret, right=of pair] (bge) {bge-small-en\\MNRL + Matryoshka};
  \node[ret, right=of bge] (idx) {chunk 512 kar.\\FAISS 1,03 jt $\times$ 384-d};
  \node[ret, right=of idx] (knn) {top-10 kosinus\\voting mayoritas};
  \draw[panah] (ocr2) -- (pair); \draw[panah] (pair) -- (bge);
  \draw[panah] (bge) -- (idx); \draw[panah] (idx) -- (knn);
\end{tikzpicture}}

  \caption{Dua pipeline Kashyap et al.\ (digambar ulang dari Gambar~1 dan~3 paper)}
  \label{fig:pipeline-kashyap}
\end{figure}

\subsection{Model per Unimodal}
\paragraph{Citra -- EfficientNetB3.} Input $300\times300$, piksel dinormalkan ke $[0,1]$,
grayscale diubah ke RGB. Keluaran 1536-d $\to$ Dropout $\to$ Dense 16 kelas. Pelatihan:
319.754 citra RVL-CDIP, 10 epoch, \emph{batch} 16, Adam lr $10^{-3}$, \emph{categorical
cross-entropy}; $\approx$22,5 jam.

\paragraph{Teks -- BERT\textsubscript{base}.} OCR memakai \emph{mesin OCR komersial}
yang tidak disebut namanya. Praproses teks: huruf kecil, tokenisasi, penyatuan spasi,
serta penghapusan angka, \emph{stopword}, karakter khusus, dan tanda baca. Tokenisasi
WordPiece diklaim mengatasi OOV. Pelatihan: 10 epoch, \emph{batch} 16, Adam lr
$5\times10^{-5}$ (adaptif), dropout 0,3; $\approx$41 jam.

\paragraph{Teks -- SCSR (bge-small-en).} Model embedding 384-d (33,3 juta parameter),
dilatih 4 epoch, \emph{batch} 16 dengan akumulasi gradien 16, AdamW lr $2\times10^{-5}$,
\emph{cosine schedule}, \emph{warmup} 10\,\%. Loss: \emph{Multiple Negatives Ranking Loss}
(MNRL) dibungkus \emph{Matryoshka loss}. Teks dipotong per 512 karakter dengan
\emph{overlap} 50 karakter, menghasilkan indeks FAISS berisi 1.030.531 vektor.

\subsection{Metode Fusi}
\textbf{Late fusion (decision-level)} berupa \emph{unweighted average ensembling}:
$\hat{y} = \arg\max_c \tfrac{1}{2}(z_v + z_t)_c$, dengan $z_v$ dan $z_t$ logit aliran
citra dan teks sebelum softmax. Tidak ada parameter fusi yang dilatih, dan bobot kedua
modalitas sama. Penulis memilih late fusion agar ekstraksi fitur terpisah dari fusi,
modular, dan galat satu modalitas tidak merambat ke modalitas lain. SCSR tidak memakai
fusi (unimodal teks).

\subsection{Metrik Evaluasi}
\begin{itemize}
  \item RVL-CDIP dan RealDoc-16: akurasi serta presisi, recall, dan F1 \emph{macro-average};
  \item Tobacco3482 (tidak seimbang): akurasi serta presisi, recall, dan F1
        \emph{weighted-average};
  \item akurasi per kelas (Tobacco3482, RealDoc-16);
  \item efisiensi: jumlah parameter, total waktu inferensi, latensi per dokumen,
        \emph{throughput};
  \item analisis kualitatif: UMAP ruang fitur, Grad-CAM (citra), \emph{integrated
        gradients} (token teks), dan matriks konfusi (lampiran).
\end{itemize}

\subsection{Hasil Keseluruhan}
Tabel~\ref{tab:kas-utama} menyalin Tabel~4 paper. Pada pendekatan~1, kedua unimodal
hampir setara di RVL-CDIP (teks 89,19 vs.\ citra 89,13), dan late fusion menaikkan
akurasi sekitar 3,9 poin ke 93,06\,\%. Di Tobacco3482, teks (90,28) sedikit di atas citra
(89,05), dan fusi mencapai 92,00\,\%. SCSR lebih rendah (80,76\,\% dan 87,89\,\%) tetapi
dianggap kompetitif untuk ukurannya. UMAP menunjukkan komplementaritas: citra memisahkan
kelas yang khas secara visual (\emph{handwritten}, \emph{advertisement}), teks memisahkan
kelas yang kaya isi (\emph{memo}, \emph{scientific publication}).

\begin{table}[H]
  \centering
  \caption{Hasil Kashyap et al.\ dibanding studi multimodal sebelumnya (Tabel~4 paper), dalam \%.
           T = teks, V = citra, T+V = fusi. F1/Rec/Prec: macro untuk RVL-CDIP, weighted untuk Tobacco3482.}
  \label{tab:kas-utama}
  \fontsize{9}{11}\selectfont
  \setlength\tabcolsep{4pt}
  \begin{tabular}{|l|c|c|c|c|c|c|c|c|c|}
    \hline
    \multirow{2}{*}{\textbf{Studi}} & \multirow{2}{*}{\textbf{Mod.}} & \multicolumn{4}{c|}{\textbf{RVL-CDIP}} & \multicolumn{4}{c|}{\textbf{Tobacco3482}} \\ \cline{3-10}
     & & Acc & F1 & Rec & Prec & Acc & F1 & Rec & Prec \\ \hline
    \multirow{3}{*}{Audebert (2019)} & T & 74,6 & -- & -- & -- & 73,8 & -- & -- & -- \\
     & V & 89,1 & -- & -- & -- & 84,5 & -- & -- & -- \\
     & T+V & 90,6 & -- & -- & -- & 87,8 & -- & -- & -- \\ \hline
    \multirow{3}{*}{Zingaro (2021)} & T & -- & -- & -- & -- & 67,8 & -- & -- & -- \\
     & V & -- & -- & -- & -- & 88,0 & -- & -- & -- \\
     & T+V & 92,7 & -- & -- & -- & 90,5 & -- & -- & -- \\ \hline
    \multirow{3}{*}{Dauphinee/Tyler (2019)} & T & 82,23 & -- & -- & -- & -- & -- & -- & -- \\
     & V & 90,24 & -- & -- & -- & -- & -- & -- & -- \\
     & T+V & 93,03 & -- & -- & -- & -- & -- & -- & -- \\ \hline
    \multirow{3}{*}{Kashyap: late fusion} & T (BERT) & 89,19 & 89,57 & 89,65 & 90,19 & 90,28 & 92,62 & 90,28 & 97,28 \\
     & V (EffNetB3) & 89,13 & 89,12 & 89,15 & 89,23 & 89,05 & 91,58 & 89,05 & 96,15 \\
     & \textbf{T+V} & \textbf{93,06} & 93,06 & 93,17 & 93,06 & \textbf{92,00} & 94,02 & 92,00 & 98,23 \\ \hline
    Kashyap: SCSR & T & 80,76 & 84,00 & 81,00 & 84,00 & 87,89 & 90,00 & 88,00 & 95,00 \\ \hline
  \end{tabular}
\end{table}

\begin{table}[H]
  \centering
  \caption{Akurasi per kelas pada Tobacco3482 (\%), Tabel~5 paper (hanya 8 dari 10 kelas dilaporkan)}
  \label{tab:kas-kelas}
  \fontsize{9}{11}\selectfont
  \setlength\tabcolsep{4pt}
  \begin{tabular}{|l|c|c|c|c|c|c|c|c|c|}
    \hline
    \textbf{Model} & \textbf{Mod.} & Letter & Form & Email & Adv & Sci-pub & News & Resume & Memo \\ \hline
    \multirow{3}{*}{Kashyap: late fusion} & T & 96,25 & 96,67 & 99,00 & 96,30 & 89,09 & 96,72 & 100,0 & 98,55 \\
     & V & 98,21 & 96,29 & 99,50 & 97,69 & 80,37 & 89,19 & 98,32 & 97,87 \\
     & T+V & 98,41 & 97,73 & 99,50 & 98,17 & 88,39 & 96,79 & 100,0 & 99,68 \\ \hline
    Kashyap: SCSR & T & 94,90 & 87,90 & 98,80 & 79,60 & 73,50 & 87,20 & 100,0 & 91,00 \\ \hline
  \end{tabular}
\end{table}

\paragraph{Efisiensi (Tabel~6--7 paper).} SCSR (33,3 juta parameter) lebih akurat di
Tobacco3482 daripada AlexNet-ELM (62 juta, 83,24\,\%) dan cabang teks Zingaro (1,8 juta,
67,8\,\%). Latensi per dokumen: late fusion 0,184~s (5,42 dok/s), SCSR 0,157~s
(6,36 dok/s); total waktu inferensi 555,27~s vs.\ 547,19~s.

\paragraph{RealDoc-16 (Tabel~8 paper).} Hanya enam kelas yang dilaporkan, akurasi
(fusi / SCSR): \emph{letter} 100/90, \emph{news article} 100/80, \emph{memo} 72,73/63,60,
\emph{form} 70/70, \emph{resume} 70/80, \emph{scientific publication} 70/90. Akurasi
keseluruhan RealDoc-16 tidak dilaporkan.

\subsection{Future Work}
\begin{itemize}
  \item Kerangka representasi yang lebih tangguh untuk dokumen \textbf{multibahasa} pada
        kondisi sumber daya rendah dan degradasi OCR berat.
  \item \emph{Parameter-efficient fine-tuning} (PEFT) untuk menekan memori pelatihan.
  \item Mengganti rata-rata tak berbobot dengan \textbf{pembobotan modalitas adaptif} atau
        fusi berbasis keyakinan (\emph{confidence-aware}), termasuk memakai model fondasi
        multimodal.
  \item Menambahkan embedding citra ringan ke basis data retrieval SCSR untuk mengatasi
        kemiripan antarkelas pada RVL-CDIP.
\end{itemize}

\subsection{Kontribusi Utama}
\begin{enumerate*}[label=(\alph*)]
  \item kerangka late fusion dua aliran (EfficientNetB3 + BERT) yang modular;
  \item SCSR, yang menggabungkan pembelajaran kontrastif terawasi dengan retrieval
        semantik sebagai pengklasifikasi tanpa lapis klasifikasi padat;
  \item dataset primer RealDoc-16 untuk uji lintas domain;
  \item analisis interpretabilitas (Grad-CAM, integrated gradients) yang memperlihatkan
        koreksi galat ganda oleh fusi.
\end{enumerate*}

\subsection{Detail Pelatihan dan Lingkungan}
Ubuntu 22.04, Intel Xeon E5-2620 v3 (12 inti), RAM 125,7~GiB, NVIDIA RTX A4000 16~GB,
CUDA 10.1, TensorFlow/Keras, Python 3.12.4. Hiperparameter tiap model sudah dirinci pada
subbagian model di atas.

\subsection{Kelebihan}
\begin{itemize}
  \item Terbit di jurnal bereputasi tahun 2026 dan \emph{open access}; memenuhi syarat
        referensi utama proyek.
  \item Membahas \emph{trade-off} akurasi--efisiensi secara eksplisit dan menawarkan dua
        titik operasi.
  \item Evaluasi pada tiga dataset termasuk data primer di luar distribusi pelatihan.
  \item Melaporkan presisi, recall, F1, dan hasil per kelas untuk setiap unimodal, serta
        analisis kualitatif (UMAP, Grad-CAM, integrated gradients, matriks konfusi).
  \item Late fusion sederhana dan mudah direplikasi; unimodal dapat dilatih dan diganti
        secara terpisah.
\end{itemize}

\subsection{Catatan Kritis}
Poin-poin berikut kami temukan saat memeriksa angka dan protokol paper.
\begin{itemize}
  \item \textbf{OCR tidak disebutkan.} Hanya disebut "mesin OCR komersial". Karena kedua
        pendekatan bergantung pada teks, tidak disebutkannya nama OCR menghalangi
        replikasi dan membuat perbandingan dengan Audebert (Tesseract) tidak setara.
  \item \textbf{Pemetaan kelas Tobacco3482 tidak dijelaskan.} Model dilatih untuk 16 kelas
        RVL-CDIP, padahal Tobacco3482 punya kelas \emph{note} dan \emph{report} yang tidak
        ada di RVL-CDIP (13,4\,\% data). Tanpa pemetaan, akurasi maksimal hanya
        $\approx$86,6\,\%, tetapi dilaporkan 92\,\%. Tabel per kelas juga hanya memuat
        8 kelas.
  \item \textbf{Potensi kebocoran data.} Tobacco3482 dan RVL-CDIP sama-sama diambil dari
        arsip Truth Tobacco Industry Documents. Paper tidak menyebut apakah dokumen
        Tobacco3482 yang juga ada di set latih RVL-CDIP sudah dibuang.
  \item \textbf{Perbandingan tidak setara protokol.} Audebert melatih pada 800 dokumen
        Tobacco3482; Kashyap melatih pada $\sim$320 ribu dokumen RVL-CDIP lalu menguji pada
        seluruh Tobacco3482. Klaim "mengungguli studi sebelumnya" di Tobacco3482 perlu
        dibaca dengan catatan ini.
  \item \textbf{Angka metrik yang tidak konsisten.}
        (i) Macro-F1 SCSR di RVL-CDIP 84,00 padahal macro-precision 84,00 dan macro-recall
        81,00. Karena F1 per kelas $\leq$ rata-rata P dan R per kelas, macro-F1 seharusnya
        $\leq$ 82,5.
        (ii) Pada set uji yang seimbang, macro-recall sama dengan akurasi, tetapi
        dilaporkan 89,65 vs.\ 89,19 (teks).
        (iii) Weighted-precision Tobacco3482 (96--98\,\%) jauh di atas recall (89--92\,\%).
        Ini wajar bila banyak prediksi jatuh ke kelas RVL-CDIP yang tidak ada di
        Tobacco3482 (dukungan nol sehingga tidak ikut dibobot), sehingga presisi tampak
        tinggi secara semu.
  \item \textbf{RealDoc-16 tidak konsisten.} Tabel~2 menyebut 160 dokumen, teks menyebut
        120. Dengan 10 dokumen per kelas, akurasi \emph{memo} 72,73\,\% (= 8/11) dan
        63,60\,\% (= 7/11) tidak mungkin. Hanya 6 dari 16 kelas yang dilaporkan dan tidak
        ada akurasi total. Jumlah data latih juga berbeda: 319.625 (Algoritma) vs.\
        319.754 (Tabel~2).
  \item \textbf{Angka Audebert dikutip keliru.} Tabel~5 paper mencantumkan akurasi per
        kelas Audebert (mis.\ \emph{letter} fusi 93,2), sedangkan Audebert hanya
        melaporkan F1 per kelas (\emph{letter} fusi 0,86) dan nilai-nilai tersebut tidak
        ditemukan di paper aslinya. Kolom citra Audebert juga ditulis "--", padahal F1
        per kelas citra dilaporkan.
  \item \textbf{Klaim efisiensi SCSR lemah.} SCSR disebut "tidak bergantung pada DL",
        padahal bge-small-en adalah transformer 33 juta parameter. Latensinya hanya
        $\sim$15\,\% lebih cepat (0,157 vs.\ 0,184~s), latensi OCR tidak dihitung, dan
        indeks FAISS 1,03 juta $\times$ 384 float32 berukuran $\approx$1,6~GB, yang
        berat untuk perangkat \emph{edge}. Selain itu, 555,27~s / 0,184~s $\approx$ 3.018
        dokumen, tidak cocok dengan 547,19~s / 0,157~s $\approx$ 3.485 dokumen jika
        keduanya diuji pada set yang sama.
  \item \textbf{Fusi tidak diablasi.} Hanya rata-rata tak berbobot yang diuji; tidak ada
        pembanding fusi berbobot, \emph{product rule}, atau fusi fitur, dan tidak ada
        oracle. Teks menyebut rata-rata logit, sedangkan Algoritma~3 menyebut probabilitas
        dan prediksi.
  \item \textbf{Tanpa variansi dan uji statistik}: hasil satu kali jalan tanpa seed, selang
        kepercayaan, atau uji signifikansi.
  \item Hal kecil: tabel menulis tahun "2025" untuk paper 2026; CUDA 10.1 tidak lazim
        dipasangkan dengan Python 3.12 dan TensorFlow versi baru; nama "Tyler et al."
        dan "Dauphinee et al." dipakai bergantian untuk rujukan yang sama.
\end{itemize}

% =====================================================================
\section{Perbandingan Kedua Paper}\label{sec:banding}

\begin{table}[H]
  \centering
  \caption{Perbandingan ringkas Audebert et al.\ (2020) dan Kashyap et al.\ (2026)}
  \label{tab:banding}
  \fontsize{9.5}{11.5}\selectfont
  \begin{tabularx}{\textwidth}{|l|>{\raggedright\arraybackslash}X|>{\raggedright\arraybackslash}X|}
    \hline
    \textbf{Aspek} & \textbf{Audebert et al.\ (2020)} & \textbf{Kashyap et al.\ (2026)} \\ \hline
    Fokus & Baseline multimodal ringan & Akurasi (fusi) dan efisiensi (SCSR) \\ \hline
    Dataset & Tobacco3482 (800 latih), RVL-CDIP & Latih RVL-CDIP; uji RVL-CDIP, Tobacco3482 penuh, RealDoc-16 \\ \hline
    OCR & Tesseract 4.0, dirilis (QS-OCR) & OCR komersial, tidak disebut \\ \hline
    Model citra & MobileNetV2, 384$\times$384 & EfficientNetB3, 300$\times$300 \\ \hline
    Model teks & FastText (beku) + CNN 1D & BERT\textsubscript{base} (fine-tune); bge-small-en (SCSR) \\ \hline
    Fusi & Fitur (konkatenasi 128+128) + MLP, end-to-end & Keputusan (rata-rata logit tak berbobot) \\ \hline
    Metrik & OA, macro-F1, F1/kelas, oracle, latensi & Acc, P/R/F1 macro/weighted, acc/kelas, parameter, latensi \\ \hline
    Unimodal terkuat & Citra jauh di atas teks & Teks $\approx$ citra (teks sedikit unggul) \\ \hline
    Fusi Tobacco3482 & 87,8\,\% (+3,3 dari citra) & 92,00\,\% (+1,72 dari teks) \\ \hline
    Fusi RVL-CDIP & 90,6\,\% (+1,5 dari citra) & 93,06\,\% (+3,87 dari teks) \\ \hline
    Reproduktibilitas & Tinggi: data teks terbuka, hiperparameter lengkap & Sedang: OCR anonim, RealDoc-16 tertutup, angka tidak konsisten \\ \hline
  \end{tabularx}
\end{table}

Ada tiga perbedaan penting. Pertama, \textbf{kualitas cabang teks menentukan nilai fusi}.
Pada Audebert, teks tertinggal 10--15 poin dari citra, sedangkan pada Kashyap teks setara
atau sedikit lebih baik dari citra. Kenaikan ini datang dari dua hal yang tidak
dipisahkan di Kashyap: OCR yang berbeda dan model teks yang lebih kuat (BERT vs.\ FastText
+ CNN). Kedua, \textbf{titik fusi}: Audebert menggabungkan fitur dan melatih fusi,
sedangkan Kashyap merata-rata keputusan tanpa pelatihan. Ketiga, \textbf{protokol
Tobacco3482} berbeda (latih 800 vs.\ \emph{transfer} dari RVL-CDIP), sehingga angka
Tobacco3482 keduanya tidak bisa dibandingkan langsung.

% =====================================================================
\section{Relevansi untuk Proyek}\label{sec:relevansi}

Fase~2 proyek kami menjadikan \textbf{mesin OCR} dan \textbf{model embedding} sebagai
variabel. Dari review ini:
\begin{itemize}
  \item Kedua paper tidak memisahkan pengaruh OCR dari pengaruh model teks. Audebert
        memakai satu OCR, Kashyap memakai OCR anonim sekaligus mengganti model teks.
        Eksperimen faktorial OCR $\times$ embedding pada proyek kami mengisi celah ini
        secara langsung.
  \item Late fusion Kashyap (rata-rata logit) murah dan dapat ditambahkan sebagai
        pembanding fusi konkatenasi Audebert, bersama oracle sebagai batas atas.
  \item Pelajaran metodologis yang kami terapkan: laporkan rerata $\pm$ simpangan baku
        lintas seed, gunakan macro-F1 untuk data tidak seimbang, nyatakan protokol split
        dan mesin OCR secara eksplisit, serta jangan membandingkan angka antar-protokol
        yang berbeda.
\end{itemize}

\section{Kesimpulan}
Audebert et al.\ memberi baseline multimodal yang sederhana, transparan, dan dapat
direplikasi. Kekuatannya adalah analisis derau OCR, oracle, dan F1 per kelas; kelemahannya
adalah satu konfigurasi OCR dan tidak ada variansi. Kashyap et al.\ memperbarui komponen
(EfficientNetB3, BERT) dan menunjukkan bahwa late fusion sederhana sudah memberi akurasi
tinggi, serta menawarkan alternatif retrieval yang ringan. Namun paper ini memiliki
sejumlah masalah pelaporan: OCR anonim, pemetaan kelas Tobacco3482 yang tidak jelas,
beberapa angka metrik yang tidak konsisten secara matematis, dan kutipan hasil Audebert
yang tidak cocok dengan sumbernya. Angka dari paper kedua sebaiknya dipakai sebagai acuan
arah, bukan target yang harus dicapai persis.


\printbibliography[heading=laporan]
\end{document}
